% --- Archon physics formalization source begin ---
% archon:physics
% archon:covers IPhO2026Problems/problem_IPhO_2026_2_A_1.lean
% archon:source-report reports/ipho_2026/problem_IPhO_2026_2_A_1.source.json
% archon:problem-id IPhO_2026_2
% archon:part-id A.1

\chapter{Physics problem IPhO\_2026\_2\_A\_1}
\label{ch:IPhO2026Problems_problem_IPhO_2026_2_A_1}

\paragraph{Problem source.}
Parallel rays strike the inside of a half-cylindrical mirror of radius R.  For
an incident ray with transverse coordinate x, let N be its number of
reflections.  The positive threshold x\_N is the largest distance from the
optical axis for which a ray undergoes at most N reflections.  Use Figures
2c--2e for the mirror and limiting-ray geometry.

Current subquestion:
Find the general expression for the threshold x\_N in terms of R and the positive integer N.

\paragraph{Current subquestion.}
Find the general expression for the threshold x\_N in terms of R and the positive integer N.

\paragraph{Recorded answer/context.}
x\_N = R*sin((2*N - 1)*pi/(4*N + 2)) = R*cos(pi/(2*N + 1)).

\paragraph{Figure/image path.}
/root/proposal\_for\_physic/science-mango/ipho\_2026\_source/image/T2\_page-2.png

\paragraph{Formalization target.}
create a compiling Lean file with sorry bodies at `IPhO2026Problems/problem\_IPhO\_2026\_2\_A\_1.lean`. The Lean declarations must preserve the physical quantities, dimensions or dimensional roles, figure labels, governing-law hypotheses, and final relation expressed by this problem.
Use Mathlib/Physlib names found through LeanExplore where available. If a physics API is missing, introduce faithful local abstractions rather than scalar placeholder aliases.

\paragraph{Iteration 002 redraft contract.}
Import the available Physlib dimensional framework and represent the mirror
radius, incident transverse coordinate, and threshold as length quantities.
Equations involving trigonometric functions may use one explicitly named SI
length projection; angles and direction components remain dimensionless.  The
positive-threshold predicate must still express attainment and maximality, and
the equal-angle reflection law must still produce the limiting-ray projection
and the \((2N+1)\)-angle closure.  Do not retain the current undocumented
``same fixed unit'' bare-real convention.

\begin{theorem}[Physics formalization target]
  \label{thm:physics:IPhO_2026_2_A_1:target}
  \lean{IPhO2026Problem2A1.positive_reflection_threshold_formula}
  \uses{def:physics:IPhO_2026_2_A_1:length_quantity, def:physics:IPhO_2026_2_A_1:si_length_value, def:physics:IPhO_2026_2_A_1:aux001, def:physics:IPhO_2026_2_A_1:aux002, def:physics:IPhO_2026_2_A_1:aux003, def:physics:IPhO_2026_2_A_1:aux004, def:physics:IPhO_2026_2_A_1:aux005, def:physics:IPhO_2026_2_A_1:aux006, def:physics:IPhO_2026_2_A_1:aux007, def:physics:IPhO_2026_2_A_1:aux008, def:physics:IPhO_2026_2_A_1:aux009, def:physics:IPhO_2026_2_A_1:aux010, def:physics:IPhO_2026_2_A_1:aux011, def:physics:IPhO_2026_2_A_1:aux012, lem:physics:IPhO_2026_2_A_1:aux013, lem:physics:IPhO_2026_2_A_1:aux014, lem:physics:IPhO_2026_2_A_1:aux015}
  IPhO 2026 Problem 2 A.1: the positive threshold for at most N reflections in the half-cylindrical mirror has the two equivalent official closed forms.
\end{theorem}
\begin{proof}
  For the positive limiting ray, repeated equal-angle reflection across the
  semicircle gives \((2N+1)\theta=\pi\).  Its transverse SI projection is
  \(R\cos\theta\), hence \(x_N=R\cos(\pi/(2N+1))\).  The complementary-angle
  identity converts this to the stated sine form; positivity of \(N\) and the
  acute limiting branch select the required threshold.
\end{proof}
\section*{Formal declaration index}
The following blocks record the typed carriers and bridge declarations used by the target.  Their dependency lists follow the declaration-level model in the covered Lean file.

\begin{definition}[Declaration LengthQuantity]
  \label{def:physics:IPhO_2026_2_A_1:length_quantity}
  \lean{IPhO2026Problem2A1.LengthQuantity}
  A physical length represented by Physlib's unit-dependent dimensionful
  carrier with length dimension.
\end{definition}
\begin{proof}
  This is the specialization of Physlib's dimensionful-quantity interface to
  the physical dimension of length.
\end{proof}

\begin{definition}[Declaration siLengthValue]
  \label{def:physics:IPhO_2026_2_A_1:si_length_value}
  \lean{IPhO2026Problem2A1.siLengthValue}
  \uses{def:physics:IPhO_2026_2_A_1:length_quantity}
  The scalar value in metres obtained by evaluating a physical length with
  Physlib's SI unit choice.
\end{definition}
\begin{proof}
  Evaluate the dimensionful length at the SI unit choice and take the
  underlying real readout.
\end{proof}

\begin{definition}[Declaration SourceFigure]
  \label{def:physics:IPhO_2026_2_A_1:aux001}
  \lean{IPhO2026Problem2A1.SourceFigure}
  The three source panels used for the multiple-reflection model: Figure 2c is the physical half-cylinder, Figure 2d is its planar cross-section, and Figure 2e is the reflection-count plot.
\end{definition}
\begin{proof}
  This is the typed definition of the stated physical carrier; its fields or defining equation provide the claimed data.
\end{proof}

\begin{definition}[Declaration AxialDirection]
  \label{def:physics:IPhO_2026_2_A_1:aux002}
  \lean{IPhO2026Problem2A1.AxialDirection}
  The orientation of a ray parallel to the optical axis in the coordinate frame of Figure 2d.
\end{definition}
\begin{proof}
  This is the typed definition of the stated physical carrier; its fields or defining equation provide the claimed data.
\end{proof}

\begin{definition}[Declaration HalfCylindricalMirror]
  \label{def:physics:IPhO_2026_2_A_1:aux003}
  \lean{IPhO2026Problem2A1.HalfCylindricalMirror}
  \uses{def:physics:IPhO_2026_2_A_1:length_quantity, def:physics:IPhO_2026_2_A_1:si_length_value}
  A half-cylindrical mirror with a positive Physlib length-valued radius.
  Coordinate formulae use only the carrier's explicitly named SI projection;
  the carrier itself, rather than a bare real, is the physical radius.
\end{definition}
\begin{proof}
  This is the typed definition of the stated physical carrier; its fields or defining equation provide the claimed data.
\end{proof}

\begin{definition}[Declaration HalfCylindricalMirror.diameter]
  \label{def:physics:IPhO_2026_2_A_1:aux004}
  \lean{IPhO2026Problem2A1.HalfCylindricalMirror.diameter}
  \uses{def:physics:IPhO_2026_2_A_1:length_quantity, def:physics:IPhO_2026_2_A_1:aux003}
  The diameter marked 2R in Figure 2c.
\end{definition}
\begin{proof}
  This is the typed definition of the stated physical carrier; its fields or defining equation provide the claimed data.
\end{proof}

\begin{definition}[Declaration OnReflectingSemicircle]
  \label{def:physics:IPhO_2026_2_A_1:aux005}
  \lean{IPhO2026Problem2A1.OnReflectingSemicircle}
  \uses{def:physics:IPhO_2026_2_A_1:length_quantity, def:physics:IPhO_2026_2_A_1:si_length_value, def:physics:IPhO_2026_2_A_1:aux003}
  The reflecting upper semicircle in the coordinate frame of Figure 2d. The optical axis is the y-axis and the origin is the circle center.
\end{definition}
\begin{proof}
  This is the typed definition of the stated physical carrier; its fields or defining equation provide the claimed data.
\end{proof}

\begin{definition}[Declaration MultipleReflectionExperiment]
  \label{def:physics:IPhO_2026_2_A_1:aux006}
  \lean{IPhO2026Problem2A1.MultipleReflectionExperiment}
  \uses{def:physics:IPhO_2026_2_A_1:length_quantity, def:physics:IPhO_2026_2_A_1:si_length_value, def:physics:IPhO_2026_2_A_1:aux002, def:physics:IPhO_2026_2_A_1:aux003}
  Typed ray data for Figures 2d--2e.  The transverse coordinate is a Physlib
  length quantity (with a named SI projection for analytic formulae), the angle
  is dimensionless, and reflectionCount records the discrete number of impacts.
\end{definition}
\begin{proof}
  This is the typed definition of the stated physical carrier; its fields or defining equation provide the claimed data.
\end{proof}

\begin{definition}[Declaration ObeysSpecularReflection]
  \label{def:physics:IPhO_2026_2_A_1:aux007}
  \lean{IPhO2026Problem2A1.ObeysSpecularReflection}
  \uses{def:physics:IPhO_2026_2_A_1:aux003, def:physics:IPhO_2026_2_A_1:aux006}
  The equal-angle law of specular reflection, stated at every actual impact. Angles are dimensionless real readouts in radians.
\end{definition}
\begin{proof}
  This is the typed definition of the stated physical carrier; its fields or defining equation provide the claimed data.
\end{proof}

\begin{definition}[Declaration IsPositiveReflectionThreshold]
  \label{def:physics:IPhO_2026_2_A_1:aux008}
  \lean{IPhO2026Problem2A1.IsPositiveReflectionThreshold}
  \uses{def:physics:IPhO_2026_2_A_1:aux003, def:physics:IPhO_2026_2_A_1:aux006}
  xN is the positive maximum transverse distance among incident rays that undergo at most N reflections. This is the threshold notion depicted in Figure 2e; it is not a closed-form definition of the requested answer.
\end{definition}
\begin{proof}
  This is the typed definition of the stated physical carrier; its fields or defining equation provide the claimed data.
\end{proof}

\begin{definition}[Declaration LimitingRayWitness]
  \label{def:physics:IPhO_2026_2_A_1:aux009}
  \lean{IPhO2026Problem2A1.LimitingRayWitness}
  \uses{def:physics:IPhO_2026_2_A_1:aux003, def:physics:IPhO_2026_2_A_1:aux006}
  A limiting ray on the positive-x branch, together with the acute polar angle of its first impact point, measured from the positive x-axis. This is distinct from the optical incidence angle relative to the surface normal.
\end{definition}
\begin{proof}
  This is the typed definition of the stated physical carrier; its fields or defining equation provide the claimed data.
\end{proof}

\begin{definition}[Declaration HalfCircleProjectionGeometry]
  \label{def:physics:IPhO_2026_2_A_1:aux010}
  \lean{IPhO2026Problem2A1.HalfCircleProjectionGeometry}
  \uses{def:physics:IPhO_2026_2_A_1:aux003, def:physics:IPhO_2026_2_A_1:aux006, def:physics:IPhO_2026_2_A_1:aux009}
  The projection read from the right triangle in the circular cross-section: the positive transverse coordinate is R cos θ.
\end{definition}
\begin{proof}
  This is the typed definition of the stated physical carrier; its fields or defining equation provide the claimed data.
\end{proof}

\begin{definition}[Declaration RepeatedReflectionClosure]
  \label{def:physics:IPhO_2026_2_A_1:aux011}
  \lean{IPhO2026Problem2A1.RepeatedReflectionClosure}
  \uses{def:physics:IPhO_2026_2_A_1:aux003, def:physics:IPhO_2026_2_A_1:aux006, def:physics:IPhO_2026_2_A_1:aux009}
  The angular closure obtained by applying equal-angle reflection repeatedly to the limiting ray in a semicircle. The 2N + 1 factor retains the endpoint and orientation information of the positive limiting branch.
\end{definition}
\begin{proof}
  This is the typed definition of the stated physical carrier; its fields or defining equation provide the claimed data.
\end{proof}

\begin{definition}[Declaration HalfCylinderReflectionLaws]
  \label{def:physics:IPhO_2026_2_A_1:aux012}
  \lean{IPhO2026Problem2A1.HalfCylinderReflectionLaws}
  \uses{def:physics:IPhO_2026_2_A_1:aux003, def:physics:IPhO_2026_2_A_1:aux006, def:physics:IPhO_2026_2_A_1:aux007, def:physics:IPhO_2026_2_A_1:aux008, def:physics:IPhO_2026_2_A_1:aux009, def:physics:IPhO_2026_2_A_1:aux010, def:physics:IPhO_2026_2_A_1:aux011}
  The governing-law interface needed from geometric optics. Besides the local equal-angle law, it exposes the limiting-ray projection and angular closure that follow from applying that law to the circular mirror.
\end{definition}
\begin{proof}
  This is the typed definition of the stated physical carrier; its fields or defining equation provide the claimed data.
\end{proof}

\begin{lemma}[Declaration limiting\_first\_impact\_angle]
  \label{lem:physics:IPhO_2026_2_A_1:aux013}
  \lean{IPhO2026Problem2A1.limiting_first_impact_angle}
  \uses{def:physics:IPhO_2026_2_A_1:aux003, def:physics:IPhO_2026_2_A_1:aux006, def:physics:IPhO_2026_2_A_1:aux009, def:physics:IPhO_2026_2_A_1:aux011}
  Algebraic bridge from the repeated-reflection closure to the unique limiting angle.
\end{lemma}
\begin{proof}
  Substitute the cited governing relations in order and use the stated positivity, orientation, and branch conditions to obtain the conclusion.
\end{proof}

\begin{lemma}[Declaration official\_answer\_angles\_complementary]
  \label{lem:physics:IPhO_2026_2_A_1:aux014}
  \lean{IPhO2026Problem2A1.official_answer_angles_complementary}
  \uses{lem:physics:IPhO_2026_2_A_1:aux013}
  The two angles occurring in the official sine and cosine answer forms are complementary.
\end{lemma}
\begin{proof}
  Substitute the cited governing relations in order and use the stated positivity, orientation, and branch conditions to obtain the conclusion.
\end{proof}

\begin{lemma}[Declaration official\_sine\_cosine\_forms\_agree]
  \label{lem:physics:IPhO_2026_2_A_1:aux015}
  \lean{IPhO2026Problem2A1.official_sine_cosine_forms_agree}
  \uses{lem:physics:IPhO_2026_2_A_1:aux013, lem:physics:IPhO_2026_2_A_1:aux014}
  Trigonometric bridge between the two official closed forms. The Mathlib carrier for the complementary-angle step is Real.sin\_pi\_div\_two\_sub.
\end{lemma}
\begin{proof}
  Substitute the cited governing relations in order and use the stated positivity, orientation, and branch conditions to obtain the conclusion.
\end{proof}
% --- Archon physics formalization source end ---
